\section{Häsordnung}

\subsection{Präambel}
Diese Häsordnung regelt das gemeinsame Auftreten, Verhalten und die Pflege der Häs innerhalb unseres Vereins. Sie gilt für alle Mitglieder und Maskengruppen gleichermaßen. \\
Die Häsordnung besteht aus einem allgemeinen Teil, der für alle Gruppen verbindlich ist, und aus gruppenspezifischen Bestimmungen, die Besonderheiten einzelner Maskengruppen regeln. Abweichungen von dieser Ordnung bedürfen der Genehmigung durch den Vorstand.

\subsection{Allgemeine Häsordnung}
\subsubsection{§ 1 Erlass, Änderung, Aufhebung und Bekanntmachung}
\begin{enumerate}
    \item Diese Häsordnung kann jederzeit durch den Vereinsausschuss erarbeitet und beschlossen werden.

    \item Änderungen werden durch den Beschluss des Vereinsausschusses wirksam und allen Mitgliedern bekannt gemacht.
\end{enumerate}

\subsubsection{§ 2 Auftreten}
\begin{enumerate}
    \item Das Häs ist bei sämtlichen fasnachtlichen Aktivitäten mit der Gruppe grundsätzlich komplett (siehe Häsbestandteile der jeweiligen Gruppe) zu tragen.

    \item Umzüge werden grundsätzlich im vollständigen Häs einschließlich Maske gelaufen. Maske sowie Handschuhe können bei Veranstaltungen, die keinen Umzug umrahmen, zu Hause gelassen werden, sofern der Vorstand vor der Veranstaltung nichts anderes festgelegt hat.

    \item Die Tragepflicht der Maske während eines Umzuges kann nur in begründeten Ausnahmefällen aufgehoben werden, beispielsweise bei Aufsichtspersonen von Kindern, aus medizinischen Gründen oder wenn es die Aktivität während des Umzuges zwingend erfordert. Über solche Ausnahmen entscheidet ausschließlich der Vorstand in verbindung mit dem Häswart.

    \item Das Tragen einer Holzmaske ist erst ab einem Alter von 14 Jahren gestattet. Mit dem Erreichen dieses Alters erhält das Mitglied eine eigene Laufnummer.

    \item Eltern oder Aufsichtspersonen, die mit Kindern teilnehmen, laufen grundsätzlich vor der Gruppe, um eine bessere Sichtbarkeit sicherzustellen und Unfälle zu vermeiden.

    \item Das komplette Häs muss zu jeder Veranstaltung ordentlich und gepflegt sein.

    \item Unter dem Häs, sowie bei außerfasnachtlichen Aktivitäten wird Vereinskleidung getragen.

    \item Vereinskleidung, das vollständige Häs sowie die Maske dürfen ausschließlich von dazu berechtigten Personen und Mitgliedern erworben und getragen werden, da diese Bestandteile wesentliche Identifikationsmerkmale des Vereins darstellen. Der Ausschuss kann im Einzelfall Sondergenehmigungen erteilen.
\end{enumerate}

\subsubsection{§ 3 Häskontrolle}
\begin{enumerate}
    \item Einmal jährlich findet vor Beginn der Kampagne eine Häskontrolle statt. Diese wird vom Häswart durchgeführt, der Termin wird in Abstimmung mit dem Vorstand festgelegt.

    \item Werden bei der Häskontrolle Mängel festgestellt, hat das Mitglied diese eigenständig und fristgerecht zu beheben. Das Häs ist anschließend zu einem mit dem Häswart vereinbarten Termin erneut vorzuzeigen und abnehmen zu lassen.

    \item Kann ein Mitglied an dem angesetzten Termin der Häskontrolle nicht teilnehmen, ist es verpflichtet, selbstständig und rechtzeitig mit dem Häswart einen Ersatztermin zu vereinbaren.

    \item Die Teilnahme an Veranstaltungen während der Kampagne ist nur mit einem vollständig geprüften und abgenommenen Häs gestattet.
\end{enumerate}

\subsubsection{§ 4 Zeitwertermittlung beim Austritt}
\begin{enumerate}
    \item Beim Austritt aus dem Verein kann der Verein das Häs gemäß § 16 Abs. 5 der Satzung zum Zeitwert zurückkaufen. Ein Anspruch auf Rückkauf besteht nicht; der Verein besitzt jedoch grundsätzlich das Vorkaufsrecht, die Entscheidung trifft der Ausschuss.

    \item Der Zeitwert des Häs wird nach einem allgemeinen Abschreibungsgrundsatz berechnet. Grundlage ist der ursprüngliche Anschaffungspreis des genehmigten Häses.

    \item Der Abschreibungszeitraum beträgt 10 Jahre. Der Wert vermindert sich jährlich um 10 \% des ursprünglichen Anschaffungspreises. Nach Ablauf von zehn Jahren gilt das Häs als vollständig abgeschrieben und hat einen Zeitwert von 0 Euro.

    \item Weist das Häs über den üblichen Gebrauch hinausgehende Schäden, erhebliche Abnutzung oder fehlende Bestandteile auf, kann der Verein den ermittelten Zeitwert angemessen mindern.

    \item Grundlage für die Berechnung ist der Zustand des Häses zum Zeitpunkt der Rückgabe sowie eine vollständige und ordnungsgemäße Häskontrolle durch den Häswart.
\end{enumerate}

\subsection{Rotbartle}
\subsubsection{§ 5 Häsbestandteile}
\begin{enumerate}
    \item Maske mit Haaren (Schaffell) und Filzhut.
    
    \item Beige Häsjacke mit Brombeerstickerei am rechten Ärmel und Flammenapplikation am linken Ärmel sowie Lederbändel in der Farbe (Natur) in diversen Ösen.

    \item Schwarzer Ledergürtel mit Hufeisen-Gürtelschnalle, zwei Hufeisen rechtsseitig sowie einer schwarzen Ledertasche linksseitig. Die Anordnung von Hufeisen und Ledertasche kann entsprechend der Händigkeit getauscht werden.

    \item Braune Lederschürze die mit einem Lederriemen über die linke Schulter zu befestigen ist.

    \item Schwarze Zimmermannshose (vorgegebenes Muster) mit weinroten Stoff und braunen Lederflicken am rechten Hosenbein.

    \item Schwarze, knöchelhohe Stiefel.

    \item Dunkelbraune oder schwarze Handschuhe.

    \item Das Wappen ist auf der rechten Brust und die Laufnummer am linken Ärmel zu tragen.
\end{enumerate}

\subsubsection{§ 6 Ergänzungsbestandteile}
\begin{enumerate}
    \item Schwarzer Becherhalter aus Leder oder schwarzer Lederbeutel für den Trinkbecher
\end{enumerate}

\subsection{Kropfschell}
\subsubsection{§ 7 Häsbestandteile}
\begin{enumerate}
    \item Maske mit Masken/Kopftuch in den Farben Orange und Blau sowie Haaren (Rossschweif).
    
    \item Blaue Häsbluse mit dem Vereinswappen auf dem linken und der Laufnummer auf dem rechten Ärmel. Die Bluse ist über dem Rock zu tragen.
    
    \item Braunes Schultertuch aus Lederflicken und Fell mit 3 Schellen.
    
    \item Brauner Ledergürtel mit Gürtelschnalle und brauner Ledertasche rechts. Die Anordnung der Ledertasche kann entsprechend der Händigkeit getauscht werden.
    
    \item Schürze orange mit blauem Stick (Rathaus und Brunnen Scharnhausen) unten links, die Schürze wird am Gürtel befestigt.
    
    \item Brauner Rock.
    
    \item Beige lange Unterhose mit gerafftem Abschluss unten.
    
    \item Strohschuhe oder festes Schuhwerk in schwarz je nach Wetterlage.
    
    \item Ringelstulpen in den Farben blau, orange und braun.
    
    \item Dunkelbraune oder schwarze Handschuhe.
\end{enumerate}

\subsubsection{§ 8 Ergänzungsbestandteile}
\begin{enumerate}
    \item Brauner Becherhalter aus Leder oder brauner Lederbeutel für den Trinkbecher
\end{enumerate}

\subsection{Glengabachdeifl}
\subsubsection{§ 9 Häsbestandteile}
\begin{enumerate}
    \item Maske in Teufelsgestalt mit rückwärtig gekrümmten Hörnern und Langhaarziegen-kopffell als „Maskentuch“.
    
    \item Häsjacke aus Langhaarziegenfell mit verdeckten Schnallen und dem Vereinswappen auf dem linken und der Laufnummer aus Leder mit roter Nummer auf dem rechten Oberarm.
    
    \item Schwarzer oder brauner Ledergürtel passend zur Fellfarbe mit 5 Rollen in unterschiedlicher Größe und daran befestigter, farblich zum Gürtel passender Ledertasche.
    
    \item Hose aus Langhaarziegenfell mit Lederhosenträgern.
    
    \item Schwarze, knöchelhohe Stiefel oder Langhaarziegenfellschuhe.
    
    \item Schwarze Handschuhe oder Langhaarziegenfellhandschuhe.
    
    \item Alle Fellbestandteile des Häs bestehen aus Langhaarziegenfell in einer zu einander pas-senden natürlich hervorgebrachten Färbung. Die Hauptfarben des Fells sind Schwarz oder Brauntöne. Weiß darf ausschließlich für Musterungen verwendet werden.
\end{enumerate}

\subsubsection{§ 10 Egänzungsbestandteile}
\begin{enumerate}
    \item Schwarzer oder brauner Lederbeutel für den Trinkbecher je nach Gürtelfarbe.
\end{enumerate}

\subsubsection{§ 11 Besondere Regelungen}
Für die Maskengruppe "Glengabachdeifl" gelten folgende abweichende oder ergänzende Regelungen:

\begin{enumerate}
    \item Ausnahme von § 2 Abs. 1 und Abs. 2 der Allgemeinen Häsordnung \\
    Bei Veranstaltungen, die keinen Umzug umrahmen, ist es den Hästrägern gestattet zusätzlich die Häsjacke (§ 9 Abs. 2) und den Ledergürtel (§9 Abs. 3) zu Hause zu lassen.
\end{enumerate}

\subsection{Inkrafttreten}
\subsubsection{§ 12 Inkrafttreten}
\begin{enumerate}
    \item Diese Häsordnung tritt nach Beschluss der Ausschusssitzung vom dd. MMMM YYYY in Kraft. \\
    Sie ersetzt alle vorherigen Versionen
\end{enumerate}